\chapter{Meta-Workflows and Asynchronous Execution}
\label{ch:meta-async}

\begin{lessonbox}{Lesson 10 --- Compose Meta-Workflows}
\textbf{Objective.} Coordinate more than one workflow and use transversal
references.

\textbf{Concept.} \texttt{DAG\_TPS} coordinates a set of workflows. A task in
one workflow can reference a task in another workflow with
\texttt{workflow://OtherWorkflow/task/path}.

\textbf{Code} (\texttt{lesson\_10.py}):
\begin{lstlisting}
from dagon import Workflow
from dagon.dag_tps import DAG_TPS
from dagon.task import DagonTask, TaskType


if __name__ == "__main__":
    observations = Workflow("Observations")
    model = Workflow("Model")

    collect = DagonTask(
        TaskType.BATCH,
        "collect",
        "mkdir -p out; echo 'forcing' > out/forcing.txt",
    )

    run_model = DagonTask(
        TaskType.BATCH,
        "run_model",
        "cat workflow://Observations/collect/out/forcing.txt > model_output.txt",
    )

    observations.add_task(collect)
    model.add_task(run_model)

    meta = DAG_TPS("Lesson10")
    meta.add_workflow(observations)
    meta.add_workflow(model)
    meta.make_dependencies()
    meta.run()
\end{lstlisting}

\textbf{Expected behavior.} The \texttt{model} workflow waits for the
\texttt{observations} workflow task that produces \texttt{forcing.txt}.

\textbf{Verify:}
\begin{lstlisting}[style=bash]
DIR=$(find /tmp -maxdepth 1 -type d -name '*-run_model' | tail -n 1)
sed -n '1,80p' "$DIR/model_output.txt"
\end{lstlisting}

If this legacy path behaves differently in a given environment, first verify
construction without execution (\texttt{print(meta.as\_json()["name"])}) and
then inspect the dependency graph with \texttt{task.prevs} and
\texttt{task.nexts}.

\textbf{Scientific practice note.} Meta-workflows are useful when a research
system has separable but connected pipelines, such as observations, model
execution, validation, and dissemination.
\end{lessonbox}

\begin{figure}[h]
\centering
\begin{tikzpicture}[
  node distance=14mm and 10mm,
  wf/.style={rectangle, draw=dagonaccent, dashed, very thick, rounded corners,
    inner sep=8pt},
  task/.style={rectangle, rounded corners, draw=dagonblue, very thick,
    fill=blue!4, minimum width=26mm, minimum height=10mm, font=\sffamily\small},
  arr/.style={-{Stealth[length=2.6mm]}, thick, dagonblue}
]
  \node[task] (collect) {collect};
  \begin{scope}[on background layer]
    \node[wf, fit=(collect), label={[font=\sffamily\bfseries\small]above:Workflow: Observations}] {};
  \end{scope}
  \node[task, right=32mm of collect] (run_model) {run\_model};
  \begin{scope}[on background layer]
    \node[wf, fit=(run_model), label={[font=\sffamily\bfseries\small]above:Workflow: Model}] {};
  \end{scope}
  \draw[arr] (collect) -- node[above,font=\sffamily\scriptsize]{\texttt{workflow://Observations/collect/...}} (run_model);
  \node[font=\sffamily\small, below=14mm of collect, xshift=32mm] {\texttt{DAG\_TPS("Lesson10")} coordinates both workflows};
\end{tikzpicture}
\caption{A meta-workflow (\texttt{DAG\_TPS}) coordinating two independently
defined \texttt{Workflow} objects, connected by a cross-workflow
\texttt{workflow://} reference.}
\label{fig:meta-workflow}
\end{figure}

\section{Discussion: Why Compose Workflows Rather Than Tasks?}

A natural question is why \texttt{DAG\_TPS} exists at all, given that every
example so far has placed all tasks inside a single \texttt{Workflow}.
The answer is organizational, not merely technical: real research systems
frequently have \emph{separable} pipelines maintained by different people or
teams --- an observations-ingestion pipeline, a model-execution pipeline, a
validation pipeline, a dissemination pipeline --- each independently
versioned and independently runnable, yet occasionally needing to consume
each other's products. \texttt{DAG\_TPS} lets each pipeline remain a
self-contained \texttt{Workflow}, while a thin coordinating layer expresses
the few cross-pipeline data dependencies that connect them.

\begin{cautionbox}
The documentation is explicit that meta-workflow code is \emph{legacy-sensitive}
and less covered by automated tests than local single-workflow behavior. For
production scientific use, verify the exact transversal scenario with
representative workflows, and preserve the verification command alongside
the experiment record.
\end{cautionbox}

\begin{lessonbox}{Lesson 11 --- Launch a Workflow Asynchronously}
\textbf{Objective.} Start a workflow without blocking the main Python thread,
observe its lifecycle, and wait for completion only when needed.

\textbf{Concept.} \texttt{Workflow.launch()} runs the same execution path as
\texttt{run()}, but in a background thread. \texttt{Workflow.wait()} joins
that thread. Event hooks let an application report task and workflow
progress without polling task status.

\textbf{Code:}
\begin{lstlisting}
from dagon import Workflow
from dagon.task import DagonTask, TaskType


workflow = Workflow("Lesson11")
workflow.add_task(DagonTask(TaskType.BATCH, "hello", "echo hello"))

workflow.on_task_start += lambda task: print("starting", task.name)
workflow.on_task_end += lambda task: print("ended", task.name, task.status.name)
workflow.on_workflow_end += lambda wf: print("workflow complete")

workflow.launch()
print("The caller is free to do other work.")
workflow.wait()
\end{lstlisting}

\textbf{Run:}
\begin{lstlisting}[style=bash]
cp dagon.ini.sample dagon.ini
python3 lesson_11.py
\end{lstlisting}

\textbf{Expected behavior.} The caller's print statement can execute while
the workflow is running; \texttt{wait()} then blocks until the task- and
workflow-end callbacks have run. Dependencies are discovered automatically by
\texttt{launch()}; calling \texttt{make\_dependencies()} explicitly remains
useful when the graph should be inspected or validated before launch.

\textbf{Verify:}
\begin{lstlisting}[style=bash]
python3 examples/async/asynchronous_launch.py
\end{lstlisting}

\textbf{Scientific practice note.} Asynchronous launch is useful for notebook
interfaces, web services, and experiment controllers that need to remain
responsive while workflows run. Event callbacks should report progress or
update state quickly; expensive work in a callback delays the workflow
thread.
\end{lessonbox}

\section{The Workflow Lifecycle Event Model}

\texttt{launch()} returns the background \texttt{threading.Thread}; calling
\texttt{launch()} while the same workflow is already running raises
\texttt{RuntimeError}. \texttt{wait(timeout=...)} returns \texttt{True} when
the workflow completed and \texttt{False} when the timeout expired; calling
\texttt{wait()} \emph{before} \texttt{launch()} is a documented no-op that
returns \texttt{True}.

Listeners are registered with \texttt{+=} (or
\texttt{workflow.add\_listener(event\_name, callback)}) and removed with
\texttt{-=} (or \texttt{remove\_listener}). Workflow-level callbacks receive
the workflow object; task-level callbacks receive the task object. Listener
exceptions are logged and do not stop the workflow --- an important
robustness property for long-running, unattended operational deployments.

\begin{table}[h]
\centering
\caption{Workflow and task lifecycle events.}
\label{tab:lifecycle-events}
\begin{tabular}{@{}p{0.40\linewidth}p{0.16\linewidth}p{0.36\linewidth}@{}}
\toprule
\textbf{Event} & \textbf{Argument} & \textbf{Fires when} \\
\midrule
\texttt{on\_workflow\_start} & workflow & Workflow execution begins. \\
\texttt{on\_workflow\_end} & workflow & Workflow execution ends. \\
\texttt{on\_task\_start} & task & A task thread begins. \\
\texttt{on\_task\_wait} & task & The task enters dependency waiting. \\
\texttt{on\_task\_staging\_in\_start} & task & Staging-in launcher preparation begins. \\
\texttt{on\_task\_staging\_in\_end} & task & Staging-in launcher preparation completes. \\
\texttt{on\_task\_execute\_start} & task & The task begins its executor phase. \\
\texttt{on\_task\_staging\_out\_start} & task & Post-execution output/reference cleanup begins. \\
\texttt{on\_task\_staging\_out\_end} & task & Post-execution output/reference cleanup completes. \\
\texttt{on\_task\_end} & task & A task thread ends, whether finished or failed. \\
\bottomrule
\end{tabular}
\end{table}

\begin{practicebox}
DAGonStar currently generates staging-in commands as part of a task
launcher and has no separate data staging-out \emph{transfer} phase; the
corresponding hooks accurately delimit the existing preparation and cleanup
phases rather than implying a transfer step that does not yet exist. Treat
the table above as a description of current behavior, in keeping with the
project's documentation-fidelity principle.
\end{practicebox}

\section{Operational Significance for Forecasting Chains}

Asynchronous launch is the mechanism that allows DAGonStar to sit behind a
notebook interface, a web service, or an automated experiment controller
without blocking it --- exactly the operational shape required by the
recurring forecasting-chain theme of this book (developed fully in
\Cref{ch:case-study,ch:synthesis}): a controller can \texttt{launch()} a
daily forecast workflow, remain responsive to status queries or new requests
while it runs, and react to \texttt{on\_workflow\_end} to trigger
downstream publication once the chain completes.

\section{Meta-Workflows and the Modularity Argument}

It is worth strengthening, with a software-engineering argument, the
organizational rationale \Cref{ch:meta-async} already gave for
\texttt{DAG\_TPS}. Software-engineering practice has long recognized a
general principle, sometimes called the \emph{single-responsibility
principle}: a unit of code (here, a \texttt{Workflow}) should have one
reason to change. An \texttt{Observations} workflow's internal structure
should change only when the observation-ingestion process itself changes
--- a new data source, a revised quality-control step --- and should be
insulated from changes to, say, the internal structure of a \texttt{Model}
workflow's simulation stages, even though the two workflows are connected
by a data dependency.

Folding both pipelines into a single, undifferentiated \texttt{Workflow}
object would violate this principle: a researcher modifying the model's
internal task decomposition would need to navigate the observation
pipeline's tasks to find the model's, and a change to either pipeline's
internal task names could silently break \texttt{workflow://} references
authored against the old names, with no namespace boundary to localize the
blast radius of the change. \texttt{DAG\_TPS} restores the namespace
boundary the cross-workflow form \texttt{workflow://<workflow-name>/...}
already implies syntactically (\Cref{ch:schema}): each constituent
\texttt{Workflow} can be developed, version-controlled, and reasoned about
largely independently, with the meta-workflow layer responsible only for
the comparatively small number of edges that genuinely cross the boundary.

\begin{practicebox}
A practical corollary: when a researcher notices that a single
\texttt{Workflow} object has grown to encompass two or more conceptually
distinct pipelines --- observation ingestion and model execution being the
canonical example used throughout this book --- that growth is itself a
signal worth acting on. Splitting the monolithic workflow into separate
\texttt{Workflow} objects, connected through \texttt{DAG\_TPS} and explicit
cross-workflow \texttt{workflow://} references, restores the modularity
single-responsibility design favors, at the documented cost
(\Cref{ch:meta-async}) of relying on a less heavily tested code path than
single-workflow execution.
\end{practicebox}

\section{Event-Driven Design and the Observer Pattern}

The lifecycle event model of Lesson~11 (\Cref{tab:lifecycle-events}) is a
direct application of what object-oriented design literature calls the
\emph{observer pattern}: a subject (here, the \texttt{Workflow} or an
individual \texttt{Task}) maintains a list of interested observers
(callbacks registered via \texttt{+=} or \texttt{add\_listener}) and
notifies each of them, without itself needing to know anything about what
those observers do, whenever a relevant state transition occurs. This
decoupling is what allows the same workflow engine core to serve radically
different operational contexts --- a one-off script that simply prints
progress to a terminal, a notebook cell that updates a progress bar widget,
or a production web service that updates a database row and triggers a
downstream notification --- without any of those contexts requiring a
modification to the engine itself; each is simply a different set of
registered observers attached to the same, unmodified event-emitting
\texttt{Workflow}.

\begin{example}
Consider a minimal operational controller built entirely from the lifecycle
events of \Cref{tab:lifecycle-events}, with no polling loop anywhere in its
implementation:
\begin{lstlisting}
status = {"state": "idle"}

def on_start(wf):
    status["state"] = "running"

def on_end(wf):
    status["state"] = "complete"
    notify_downstream_service(wf.name)

workflow.on_workflow_start += on_start
workflow.on_workflow_end += on_end
workflow.launch()
\end{lstlisting}
This controller's \texttt{status} dictionary --- which might back, for
instance, an HTTP status endpoint in a small operational service --- stays
correctly synchronized with the workflow's actual state purely through
event notification, with no thread of the controller's own ever needing to
ask ``is it done yet?''
\end{example}

\section{Failure Isolation in the Listener Model}

\Cref{ch:meta-async} noted, in passing, that listener exceptions are logged
and do not stop the workflow. This design choice deserves a fuller
discussion, because it represents a deliberate trade-off between two
competing goals --- \emph{robustness} of the core workflow execution, and
\emph{visibility} of bugs in observer code --- and a workflow author should
understand which side of that trade-off the engine has chosen before relying
on lifecycle callbacks for anything operationally critical.

The chosen design favors robustness: a buggy or exception-raising callback
--- for instance, a progress-reporting function that crashes on an
unexpected character encoding --- cannot, by design, cause an otherwise
healthy scientific workflow to fail. This is generally the right default for
a workflow engine, since the scientific computation a workflow exists to
perform should not be held hostage to the correctness of incidental
monitoring or logging code layered on top of it.

\begin{cautionbox}
The cost of this robustness is that a callback's silent failure is, by
construction, not visible through the workflow's own \texttt{Status}
(\Cref{sec:status-model}) or return value --- only through whatever logging
configuration (\Cref{ch:configuration}) is in place to surface logged
exceptions. A workflow author who depends on a lifecycle callback for an
operationally important side effect --- triggering a downstream publication
step, say, as the synthesis chapter's forecasting-chain pattern
(\Cref{ch:synthesis}) will do --- should verify, separately from the
workflow's own success/failure status, that the relevant callback actually
executed without error, rather than assuming that workflow success implies
callback success.
\end{cautionbox}

\section{Chapter Summary}

\texttt{DAG\_TPS} composes independently defined workflows through the same
\texttt{workflow://} schema used within a single workflow, at the cost of
reduced automated-test coverage that should be respected operationally.
\texttt{launch()}/\texttt{wait()} and the lifecycle event table give a
workflow author non-blocking execution and fine-grained progress reporting
without polling. This chapter additionally grounded the meta-workflow design
in the single-responsibility principle of software engineering, framed the
lifecycle event model as an instance of the classical observer pattern with
a worked, polling-free controller example, and examined the deliberate
robustness-versus-visibility trade-off behind the engine's choice to isolate
listener exceptions from workflow execution. Together, these two lessons
complete the engine's \emph{composition and orchestration} layer;
\Cref{ch:llm-tasks} through \Cref{ch:native-tasks} now turn to the engine's
most distinctive contribution: first-class AI and web-service tasks.

\begin{exercise}
Design, in prose and pseudocode, a lifecycle-event-based mechanism that
would let a workflow author detect, \emph{without} consulting external
logs, whether an \texttt{on\_workflow\_end} callback raised an exception
during a \texttt{launch()}/\texttt{wait()} execution. (Hint: consider what a
callback can do, within itself, to communicate failure back to the code that
called \texttt{wait()}, given that the callback's own exceptions are
swallowed by the engine.)
\end{exercise}
